\section{Conclusion}

We presented SAQT, a structure-aware query learning method for vertical ancient Chinese text recognition. The method uses character boxes as training-time localization supervision for query learning rather than as final inference outputs. When the target charset differs from the reference charset, charset-aware classifier adaptation inherits shared-character classifier parameters and initializes rows for newly introduced target characters. The adapted model is then trained with line-level CTC, and direct decoding can be accompanied by a fixed development-set blank/nonblank calibration protocol at inference.

On MTHv2 and HDRC, SAQT achieves higher AR/CR than the vertical-input-adapted scene text recognition baselines considered in this paper. The HDRC random-head control supports charset-aware initialization over random target classifier initialization, and the MTHv2 query-budget ablation provides internal evidence for the length-aware query constraint. The comparison between classifier-head reconstruction and full-model training shows that updating only the output classifier is insufficient for recognition adaptation. The expected-count auxiliary term gives a small gain on MTHv2 but does not yet provide stable evidence on HDRC, so we treat it as an optional lightweight module rather than a universal core component. The no-localization CTC control further suggests that, under the current query-based architecture and training budget, line-level transcription alone is insufficient to form a usable query-to-sequence representation. Future work will study stronger charset adaptation and character-localization information learning strategies that reduce dependence on dense character-box annotations.
